\documentclass[12pt,a4paper]{article}
\usepackage[utf8]{inputenc}
\usepackage[T2A]{fontenc}
\usepackage[russian]{babel}
\usepackage{geometry}
\usepackage{longtable}
\geometry{margin=2cm}

\title{Отчёт по реализации программного средства для задачи внедрения СЗИ}
\author{}
\date{}

\begin{document}
\maketitle

\section*{Назначение реализации}
Реализовано прикладное веб-приложение для решения задачи внедрения средств защиты информации для групп информационных активов с нечёткими параметрами. Отчёт содержит только информацию о реализации программного средства.

\section*{Архитектура решения}
Реализация разделена на две части:
\begin{itemize}
    \item клиентская часть на React, отвечающая за пользовательский интерфейс, табличный ввод данных и отображение итогового результата;
    \item серверная часть на Python, выполняющая дефаззификацию нечётких параметров, сортировку данных и решение оптимизационной задачи.
\end{itemize}

\section*{Использованные инструменты}
\begin{longtable}{|p{4cm}|p{10cm}|}
\hline
\textbf{Инструмент} & \textbf{Назначение} \\\hline
Python 3 & Серверная логика, дефаззификация, решение оптимизационной задачи, локальный launcher. \\\hline
React 19 & Построение интерфейса пользователя. \\\hline
Vite & Сборка и локальная разработка клиентской части. \\\hline
Flask & HTTP API для связи интерфейса с вычислительным модулем на Python. \\\hline
Waitress & Production-ready WSGI-сервер для локального запуска приложения. \\\hline
PyInstaller & Упаковка launcher и собранного интерфейса в исполняемый файл. \\\hline
Git & Контроль версий проекта. \\\hline
\end{longtable}

\section*{Использованные библиотеки}
\begin{longtable}{|p{4cm}|p{10cm}|}
\hline
\textbf{Библиотека} & \textbf{Роль в проекте} \\\hline
NumPy & Работа с векторами, матрицами и операциями дефаззификации. \\\hline
SciPy & Решение вспомогательных булевых задач линейной оптимизации через \texttt{scipy.optimize.milp}. \\\hline
Flask & Реализация REST API для вызова вычислений из веб-интерфейса. \\\hline
Waitress & Запуск сервера без использования встроенного dev-server Flask в режиме конечного пользователя. \\\hline
React DOM & Отрисовка клиентского приложения в браузере. \\\hline
\end{longtable}

\section*{Реализованные функции интерфейса}
\begin{itemize}
    \item горизонтальная компоновка с отдельными ячейками для всех значений матриц и векторов;
    \item ввод треугольных нечётких параметров \(\tilde{c}_j\) и \(\tilde{d}_j\) в формате троек чисел;
    \item три способа задания параметра \(\lambda\): ползунок, числовое поле, набор пресетов;
    \item выбор метода дефаззификации из списка;
    \item отображение только итогового решения без промежуточных шагов.
\end{itemize}

\section*{Реализованные функции вычислительного модуля}
\begin{itemize}
    \item проверка корректности входных данных;
    \item дефаззификация треугольных нечётких чисел несколькими методами;
    \item автоматическая сортировка параметров по невозрастанию \(d_j\);
    \item решение семейства вспомогательных задач о ранце библиотечным MILP-решателем;
    \item восстановление оптимального булева вектора в исходной нумерации СЗИ;
    \item возврат итоговых значений \(c\), \(d\), \(F_1\), \(F_2\), \(F\) и оптимального решения \(x^*\).
\end{itemize}

\section*{Средства запуска и поставки}
Для удобства пользователя реализован файл \texttt{launcher.py}, который запускает локальный сервер и автоматически открывает приложение в браузере. Для сборки автономного исполняемого файла предусмотрен отдельный сценарий упаковки через PyInstaller.

\end{document}
